%% ECON 282E -- Session 5, Deck B4: Finite Elements
%% Source: Quant_Macro Lec.9 frames 88-95 (L2206-2410).  Follows
%% Fernandez-Villaverde & Guerron-Quintana, Solution Methods IV pp.32-38;
%% rewritten here and cited.  Source defects are recorded in slides/SOURCES.md,
%% not on the slides.
%% Numbers from tools/figures/s05_numbers.json, key "fem".

%% ECON 282E -- Foundations of Macroeconomics (UC Riverside, Fall 2026)
%% Shared Beamer preamble for every deck in slides/S01 ... slides/S10.
%%
%% Usage (a deck lives in slides/S0X/ and is compiled from that folder):
%%   %% ECON 282E -- Foundations of Macroeconomics (UC Riverside, Fall 2026)
%% Shared Beamer preamble for every deck in slides/S01 ... slides/S10.
%%
%% Usage (a deck lives in slides/S0X/ and is compiled from that folder):
%%   %% ECON 282E -- Foundations of Macroeconomics (UC Riverside, Fall 2026)
%% Shared Beamer preamble for every deck in slides/S01 ... slides/S10.
%%
%% Usage (a deck lives in slides/S0X/ and is compiled from that folder):
%%   \input{../preamble/course_preamble.tex}
%%   \decktitle{1}{A1}{Who I Am \& Why This Course}{September 24, 2026}
%%   \begin{document}
%%   \makedecktitle
%%   ...
%%
%% Adapted from Courses/AI-ECON-2026/source/shared_preamble.tex (Metropolis theme,
%% concept/caution/example/takeaway boxes, \sectiondivider). Changes for this course:
%%   * course title block (\decktitle / \makedecktitle) instead of \makebeamertitle
%%   * \zh{...} typesets Chinese (book titles) under pdflatex via CJKutf8
%%   * researchbox for the "Research in the wild" frames
%%   * section pages off; TikZ libraries and the course map loaded here
%% Build with pdflatex only (two passes); tools/build_slides.py does this for you.

\documentclass[10pt,english,aspectratio=169]{beamer}

\usepackage[T1]{fontenc}
\usepackage{textcomp}
\usepackage[utf8]{inputenc}
\usepackage{babel}
\usepackage{CJKutf8}

\usepackage{amsmath, amssymb, amsthm, graphicx}
\usepackage{booktabs, multirow, tabularx}
\usepackage{tikz}
\usetikzlibrary{positioning,arrows.meta,shapes,calc,fit,backgrounds}
\usepackage{pgfplots}
\pgfplotsset{compat=1.16}
\usepackage[most]{tcolorbox}
\tcbuselibrary{skins, breakable}
\usepackage{listings}
\usepackage{xcolor}

\lstset{
  basicstyle=\ttfamily\small,
  keywordstyle=\color{blue!80!black}\bfseries,
  commentstyle=\color{green!50!black}\itshape,
  stringstyle=\color{orange!80!black},
  backgroundcolor=\color{gray!8},
  frame=single,
  rulecolor=\color{gray!40},
  breaklines=true,
  showstringspaces=false,
  numbers=none,
}

\usetheme{metropolis}
\usefonttheme{professionalfonts}
\metroset{sectionpage=none}

\definecolor{LightGray}{RGB}{230,230,230}
\definecolor{RedTitle}{RGB}{0,0,200}
\definecolor{VeryLightGray}{RGB}{245,245,245}
\definecolor{DarkText}{RGB}{0,0,0}
\definecolor{UCRBlue}{RGB}{0,45,106}
\definecolor{UCRGold}{RGB}{241,171,0}

% Stage colours of the course arc (used by the course map and elsewhere)
\colorlet{StageTrunk}{blue!12}
\colorlet{StageBranch}{cyan!14}
\colorlet{StageMachine}{gray!18}
\colorlet{StageWall}{red!14}
\colorlet{StageTools}{orange!20}
\colorlet{StageData}{green!16}
\colorlet{StageAgents}{violet!14}

\hypersetup{bookmarks=false,breaklinks=true,pdfborder={0 0 0},colorlinks=true,
  linkcolor=DarkText,urlcolor=RedTitle}

\setbeamercolor{background canvas}{bg=white}
\setbeamercolor{title}{fg=RedTitle}
\setbeamercolor{frametitle}{bg=VeryLightGray,fg=DarkText}
\setbeamercolor{normal text}{fg=DarkText}
\setbeamercolor{alerted text}{fg=RedTitle}
\setbeamersize{text margin left=5mm, text margin right=5mm}
\addtobeamertemplate{frametitle}{}{\vspace{-0.5em}}

\setbeamertemplate{navigation symbols}{}
\setbeamertemplate{footline}{%
  \hspace{5mm}{\usebeamerfont{page number in head/foot}\color{black!45}%
  ECON 282E \,$\cdot$\, \insertshortdeck}\hfill%
  \usebeamercolor[fg]{page number in head/foot}%
  \usebeamerfont{page number in head/foot}%
  \insertframenumber\,/\,\inserttotalframenumber\kern1em\vskip2pt%
}

% ---------------------------------------------------------------------------
% Chinese text under pdflatex (book titles etc.)
\newcommand{\zh}[1]{\begin{CJK*}{UTF8}{gbsn}#1\end{CJK*}}

% ---------------------------------------------------------------------------
% Deck title block
%   \decktitle{session}{deck id}{title}{date}
\newcommand{\insertshortdeck}{}
\newcommand{\decksession}{}
\newcommand{\deckid}{}
\newcommand{\decktitletext}{}
\newcommand{\deckdate}{}
\newcommand{\decktitle}[4]{%
  \renewcommand{\decksession}{#1}%
  \renewcommand{\deckid}{#2}%
  \renewcommand{\decktitletext}{#3}%
  \renewcommand{\deckdate}{#4}%
  \renewcommand{\insertshortdeck}{Session #1 \,$\cdot$\, Deck #1.#2}%
  \title{#3}%
  \author{Zhigang Feng}%
  \date{#4}%
}
\newcommand{\makedecktitle}{%
  \begin{frame}[plain,noframenumbering]
    \begin{tikzpicture}[remember picture,overlay]
      \fill[UCRBlue] (current page.north west) rectangle ([yshift=-0.55cm]current page.north east);
      \fill[UCRGold] ([yshift=-0.55cm]current page.north west) rectangle ([yshift=-0.62cm]current page.north east);
      \node[anchor=west, text=white, font=\small\bfseries] at ([xshift=5mm,yshift=-0.275cm]current page.north west)
        {ECON 282E \,$\cdot$\, Foundations of Macroeconomics};
      \node[anchor=east, text=white, font=\small] at ([xshift=-5mm,yshift=-0.275cm]current page.north east)
        {UC Riverside \,$\cdot$\, Fall 2026};
    \end{tikzpicture}
    \vspace*{1.1cm}
    {\usebeamerfont{title}\color{black!55}\large Session \decksession\ \,$\cdot$\, Deck \decksession.\deckid\par}
    \vspace{0.25cm}
    {\usebeamerfont{title}\color{RedTitle}\LARGE\bfseries \decktitletext\par}
    \vspace{0.7cm}
    {\large Zhigang Feng\par}
    \vspace{0.15cm}
    {\color{black!65}\deckdate\par}
    \vfill
    {\footnotesize\itshape\color{black!55}Build it on paper $\;\to\;$ solve it on the machine $\;\to\;$ push it to the frontier with AI\par}
  \end{frame}}

% ---------------------------------------------------------------------------
% Section divider (plain frame, not counted)
\newcommand{\sectiondivider}[2]{%
\begin{frame}[plain,noframenumbering]
\begin{center}
\vfill
{\Large\textbf{#1}\par}
\vspace{0.35cm}
{\normalsize #2\par}
\vfill
\end{center}
\end{frame}
}

% ---------------------------------------------------------------------------
% Boxes
\newtcolorbox{conceptbox}[1][]{colback=blue!3, colframe=RedTitle!55, boxrule=0.35mm,
  arc=1mm, left=2mm, right=2mm, top=1mm, bottom=1mm, #1}
\newtcolorbox{cautionbox}[1][]{colback=red!3, colframe=red!50!black, boxrule=0.35mm,
  arc=1mm, left=2mm, right=2mm, top=1mm, bottom=1mm, #1}
\newtcolorbox{examplebox}[1][]{colback=green!3, colframe=green!45!black, boxrule=0.35mm,
  arc=1mm, left=2mm, right=2mm, top=1mm, bottom=1mm, #1}
\newtcolorbox{takeawaybox}[1][]{colback=VeryLightGray, colframe=RedTitle!55, boxrule=0.35mm,
  arc=1mm, left=2mm, right=2mm, top=1mm, bottom=1mm, #1}
\newtcolorbox{researchbox}[1][]{enhanced, colback=violet!4, colframe=violet!60!black,
  boxrule=0.35mm, arc=1mm, left=2mm, right=2mm, top=1mm, bottom=1mm,
  fonttitle=\bfseries\small, title={Research in the wild}, #1}

\newcommand{\compacttables}{\renewcommand{\arraystretch}{1.15}}
\newcommand{\src}[1]{{\tiny\color{black!55}#1}}

% ---------------------------------------------------------------------------
\input{../preamble/course_map.tex}

%%   \decktitle{1}{A1}{Who I Am \& Why This Course}{September 24, 2026}
%%   \begin{document}
%%   \makedecktitle
%%   ...
%%
%% Adapted from Courses/AI-ECON-2026/source/shared_preamble.tex (Metropolis theme,
%% concept/caution/example/takeaway boxes, \sectiondivider). Changes for this course:
%%   * course title block (\decktitle / \makedecktitle) instead of \makebeamertitle
%%   * \zh{...} typesets Chinese (book titles) under pdflatex via CJKutf8
%%   * researchbox for the "Research in the wild" frames
%%   * section pages off; TikZ libraries and the course map loaded here
%% Build with pdflatex only (two passes); tools/build_slides.py does this for you.

\documentclass[10pt,english,aspectratio=169]{beamer}

\usepackage[T1]{fontenc}
\usepackage{textcomp}
\usepackage[utf8]{inputenc}
\usepackage{babel}
\usepackage{CJKutf8}

\usepackage{amsmath, amssymb, amsthm, graphicx}
\usepackage{booktabs, multirow, tabularx}
\usepackage{tikz}
\usetikzlibrary{positioning,arrows.meta,shapes,calc,fit,backgrounds}
\usepackage{pgfplots}
\pgfplotsset{compat=1.16}
\usepackage[most]{tcolorbox}
\tcbuselibrary{skins, breakable}
\usepackage{listings}
\usepackage{xcolor}

\lstset{
  basicstyle=\ttfamily\small,
  keywordstyle=\color{blue!80!black}\bfseries,
  commentstyle=\color{green!50!black}\itshape,
  stringstyle=\color{orange!80!black},
  backgroundcolor=\color{gray!8},
  frame=single,
  rulecolor=\color{gray!40},
  breaklines=true,
  showstringspaces=false,
  numbers=none,
}

\usetheme{metropolis}
\usefonttheme{professionalfonts}
\metroset{sectionpage=none}

\definecolor{LightGray}{RGB}{230,230,230}
\definecolor{RedTitle}{RGB}{0,0,200}
\definecolor{VeryLightGray}{RGB}{245,245,245}
\definecolor{DarkText}{RGB}{0,0,0}
\definecolor{UCRBlue}{RGB}{0,45,106}
\definecolor{UCRGold}{RGB}{241,171,0}

% Stage colours of the course arc (used by the course map and elsewhere)
\colorlet{StageTrunk}{blue!12}
\colorlet{StageBranch}{cyan!14}
\colorlet{StageMachine}{gray!18}
\colorlet{StageWall}{red!14}
\colorlet{StageTools}{orange!20}
\colorlet{StageData}{green!16}
\colorlet{StageAgents}{violet!14}

\hypersetup{bookmarks=false,breaklinks=true,pdfborder={0 0 0},colorlinks=true,
  linkcolor=DarkText,urlcolor=RedTitle}

\setbeamercolor{background canvas}{bg=white}
\setbeamercolor{title}{fg=RedTitle}
\setbeamercolor{frametitle}{bg=VeryLightGray,fg=DarkText}
\setbeamercolor{normal text}{fg=DarkText}
\setbeamercolor{alerted text}{fg=RedTitle}
\setbeamersize{text margin left=5mm, text margin right=5mm}
\addtobeamertemplate{frametitle}{}{\vspace{-0.5em}}

\setbeamertemplate{navigation symbols}{}
\setbeamertemplate{footline}{%
  \hspace{5mm}{\usebeamerfont{page number in head/foot}\color{black!45}%
  ECON 282E \,$\cdot$\, \insertshortdeck}\hfill%
  \usebeamercolor[fg]{page number in head/foot}%
  \usebeamerfont{page number in head/foot}%
  \insertframenumber\,/\,\inserttotalframenumber\kern1em\vskip2pt%
}

% ---------------------------------------------------------------------------
% Chinese text under pdflatex (book titles etc.)
\newcommand{\zh}[1]{\begin{CJK*}{UTF8}{gbsn}#1\end{CJK*}}

% ---------------------------------------------------------------------------
% Deck title block
%   \decktitle{session}{deck id}{title}{date}
\newcommand{\insertshortdeck}{}
\newcommand{\decksession}{}
\newcommand{\deckid}{}
\newcommand{\decktitletext}{}
\newcommand{\deckdate}{}
\newcommand{\decktitle}[4]{%
  \renewcommand{\decksession}{#1}%
  \renewcommand{\deckid}{#2}%
  \renewcommand{\decktitletext}{#3}%
  \renewcommand{\deckdate}{#4}%
  \renewcommand{\insertshortdeck}{Session #1 \,$\cdot$\, Deck #1.#2}%
  \title{#3}%
  \author{Zhigang Feng}%
  \date{#4}%
}
\newcommand{\makedecktitle}{%
  \begin{frame}[plain,noframenumbering]
    \begin{tikzpicture}[remember picture,overlay]
      \fill[UCRBlue] (current page.north west) rectangle ([yshift=-0.55cm]current page.north east);
      \fill[UCRGold] ([yshift=-0.55cm]current page.north west) rectangle ([yshift=-0.62cm]current page.north east);
      \node[anchor=west, text=white, font=\small\bfseries] at ([xshift=5mm,yshift=-0.275cm]current page.north west)
        {ECON 282E \,$\cdot$\, Foundations of Macroeconomics};
      \node[anchor=east, text=white, font=\small] at ([xshift=-5mm,yshift=-0.275cm]current page.north east)
        {UC Riverside \,$\cdot$\, Fall 2026};
    \end{tikzpicture}
    \vspace*{1.1cm}
    {\usebeamerfont{title}\color{black!55}\large Session \decksession\ \,$\cdot$\, Deck \decksession.\deckid\par}
    \vspace{0.25cm}
    {\usebeamerfont{title}\color{RedTitle}\LARGE\bfseries \decktitletext\par}
    \vspace{0.7cm}
    {\large Zhigang Feng\par}
    \vspace{0.15cm}
    {\color{black!65}\deckdate\par}
    \vfill
    {\footnotesize\itshape\color{black!55}Build it on paper $\;\to\;$ solve it on the machine $\;\to\;$ push it to the frontier with AI\par}
  \end{frame}}

% ---------------------------------------------------------------------------
% Section divider (plain frame, not counted)
\newcommand{\sectiondivider}[2]{%
\begin{frame}[plain,noframenumbering]
\begin{center}
\vfill
{\Large\textbf{#1}\par}
\vspace{0.35cm}
{\normalsize #2\par}
\vfill
\end{center}
\end{frame}
}

% ---------------------------------------------------------------------------
% Boxes
\newtcolorbox{conceptbox}[1][]{colback=blue!3, colframe=RedTitle!55, boxrule=0.35mm,
  arc=1mm, left=2mm, right=2mm, top=1mm, bottom=1mm, #1}
\newtcolorbox{cautionbox}[1][]{colback=red!3, colframe=red!50!black, boxrule=0.35mm,
  arc=1mm, left=2mm, right=2mm, top=1mm, bottom=1mm, #1}
\newtcolorbox{examplebox}[1][]{colback=green!3, colframe=green!45!black, boxrule=0.35mm,
  arc=1mm, left=2mm, right=2mm, top=1mm, bottom=1mm, #1}
\newtcolorbox{takeawaybox}[1][]{colback=VeryLightGray, colframe=RedTitle!55, boxrule=0.35mm,
  arc=1mm, left=2mm, right=2mm, top=1mm, bottom=1mm, #1}
\newtcolorbox{researchbox}[1][]{enhanced, colback=violet!4, colframe=violet!60!black,
  boxrule=0.35mm, arc=1mm, left=2mm, right=2mm, top=1mm, bottom=1mm,
  fonttitle=\bfseries\small, title={Research in the wild}, #1}

\newcommand{\compacttables}{\renewcommand{\arraystretch}{1.15}}
\newcommand{\src}[1]{{\tiny\color{black!55}#1}}

% ---------------------------------------------------------------------------
%% Course map: the recurring "you are here" slide for ECON 282E.
%% Loaded by course_preamble.tex. Pattern follows AI-ECON-2026 lec01_map.tex, but the map
%% shows the whole ten-session course rather than one lecture.
%%
%%   \CourseMap{s}{label}   s = 1..10 highlights session s and prints "you are here: label"
%%                          (e.g. \CourseMap{1}{1.A2}); earlier sessions get a check mark.
%%   \CourseMap{0}{}        full overview, nothing highlighted.
%%
%% Note: TikZ option lists are comma-split before TeX conditionals run, so every \ifnum
%% below sits outside [...] option lists.

\newcommand{\CourseMap}[2]{%
\begin{center}
\begin{tikzpicture}[
    sess/.style={rectangle, rounded corners=2pt, draw=black!45, minimum width=1.34cm,
        minimum height=1.12cm, text width=1.26cm, align=center, inner sep=1pt},
    stage/.style={font=\tiny\bfseries, text=black!70},
]
  \foreach \i/\d/\t/\c in {%
      1/Sep 24/Intro \\ Growth/StageTrunk,
      2/Oct 1/HA $\cdot$ OLG \\ Policy/StageBranch,
      3/Oct 8/Num.\ DP \\ Python/StageMachine,
      4/Oct 15/Numerics \\ I/StageMachine,
      5/Oct 22/Numerics \\ II/StageMachine,
      6/Oct 29/The wall \\ \textbf{Midterm}/StageWall,
      7/Nov 5/ML \\ Deep nets/StageTools,
      8/Nov 12/RL \\ HA $+$ DL/StageTools,
      9/Nov 19/Text \\ LLMs/StageData,
      10/Dec 3/Agents \\ \textbf{Projects}/StageAgents} {
    \node[sess, fill=\c] (n\i) at ({1.46*(\i-1)}, 0)
      {{\scriptsize\bfseries S\i}\\[-1pt]{\tiny\color{black!60}\d}\\[1pt]{\tiny \t}};
    \ifnum\i<#1
      \node[font=\scriptsize, text=green!45!black] at ([xshift=-2.5mm,yshift=-2.2mm]n\i.north east) {$\checkmark$};
    \fi
    \ifnum\i=#1
      \draw[RedTitle, very thick, rounded corners=2pt]
        ([xshift=-0.6mm,yshift=-0.6mm]n\i.south west) rectangle ([xshift=0.6mm,yshift=0.6mm]n\i.north east);
      % keep the label inside the picture at both ends of the row
      \ifnum\i<3
        \node[font=\scriptsize\bfseries, text=RedTitle, anchor=south west] at ([yshift=1.2mm]n\i.north west)
          {$\blacktriangledown$\,you are here: #2};
      \else\ifnum\i>8
        \node[font=\scriptsize\bfseries, text=RedTitle, anchor=south east] at ([yshift=1.2mm]n\i.north east)
          {you are here: #2\,$\blacktriangledown$};
      \else
        \node[font=\scriptsize\bfseries, text=RedTitle, anchor=south] at ([yshift=1.2mm]n\i.north)
          {you are here: #2\,$\blacktriangledown$};
      \fi\fi
    \fi
  }
  % stage band under the sessions
  \foreach \a/\b/\lab/\c in {1/1/trunk/StageTrunk, 2/2/branches/StageBranch,
      3/5/the machine $\cdot$ classical methods/StageMachine, 6/6/the wall/StageWall,
      7/8/new tools/StageTools, 9/9/new data/StageData, 10/10/colleagues/StageAgents} {
    \fill[\c] ([yshift=-1.5mm]n\a.south west) rectangle ([yshift=-4.6mm]n\b.south east);
    \path ([yshift=-1.5mm]n\a.south west) -- ([yshift=-4.6mm]n\b.south east)
      node[midway, stage] {\lab};
  }
  \node[font=\footnotesize\itshape, text=black!70] at ([yshift=-9.5mm]$(n1.south)!0.5!(n10.south)$)
    {Build it on paper $\;\to\;$ solve it on the machine $\;\to\;$ push it to the frontier with AI};
\end{tikzpicture}
\end{center}
}


%%   \decktitle{1}{A1}{Who I Am \& Why This Course}{September 24, 2026}
%%   \begin{document}
%%   \makedecktitle
%%   ...
%%
%% Adapted from Courses/AI-ECON-2026/source/shared_preamble.tex (Metropolis theme,
%% concept/caution/example/takeaway boxes, \sectiondivider). Changes for this course:
%%   * course title block (\decktitle / \makedecktitle) instead of \makebeamertitle
%%   * \zh{...} typesets Chinese (book titles) under pdflatex via CJKutf8
%%   * researchbox for the "Research in the wild" frames
%%   * section pages off; TikZ libraries and the course map loaded here
%% Build with pdflatex only (two passes); tools/build_slides.py does this for you.

\documentclass[10pt,english,aspectratio=169]{beamer}

\usepackage[T1]{fontenc}
\usepackage{textcomp}
\usepackage[utf8]{inputenc}
\usepackage{babel}
\usepackage{CJKutf8}

\usepackage{amsmath, amssymb, amsthm, graphicx}
\usepackage{booktabs, multirow, tabularx}
\usepackage{tikz}
\usetikzlibrary{positioning,arrows.meta,shapes,calc,fit,backgrounds}
\usepackage{pgfplots}
\pgfplotsset{compat=1.16}
\usepackage[most]{tcolorbox}
\tcbuselibrary{skins, breakable}
\usepackage{listings}
\usepackage{xcolor}

\lstset{
  basicstyle=\ttfamily\small,
  keywordstyle=\color{blue!80!black}\bfseries,
  commentstyle=\color{green!50!black}\itshape,
  stringstyle=\color{orange!80!black},
  backgroundcolor=\color{gray!8},
  frame=single,
  rulecolor=\color{gray!40},
  breaklines=true,
  showstringspaces=false,
  numbers=none,
}

\usetheme{metropolis}
\usefonttheme{professionalfonts}
\metroset{sectionpage=none}

\definecolor{LightGray}{RGB}{230,230,230}
\definecolor{RedTitle}{RGB}{0,0,200}
\definecolor{VeryLightGray}{RGB}{245,245,245}
\definecolor{DarkText}{RGB}{0,0,0}
\definecolor{UCRBlue}{RGB}{0,45,106}
\definecolor{UCRGold}{RGB}{241,171,0}

% Stage colours of the course arc (used by the course map and elsewhere)
\colorlet{StageTrunk}{blue!12}
\colorlet{StageBranch}{cyan!14}
\colorlet{StageMachine}{gray!18}
\colorlet{StageWall}{red!14}
\colorlet{StageTools}{orange!20}
\colorlet{StageData}{green!16}
\colorlet{StageAgents}{violet!14}

\hypersetup{bookmarks=false,breaklinks=true,pdfborder={0 0 0},colorlinks=true,
  linkcolor=DarkText,urlcolor=RedTitle}

\setbeamercolor{background canvas}{bg=white}
\setbeamercolor{title}{fg=RedTitle}
\setbeamercolor{frametitle}{bg=VeryLightGray,fg=DarkText}
\setbeamercolor{normal text}{fg=DarkText}
\setbeamercolor{alerted text}{fg=RedTitle}
\setbeamersize{text margin left=5mm, text margin right=5mm}
\addtobeamertemplate{frametitle}{}{\vspace{-0.5em}}

\setbeamertemplate{navigation symbols}{}
\setbeamertemplate{footline}{%
  \hspace{5mm}{\usebeamerfont{page number in head/foot}\color{black!45}%
  ECON 282E \,$\cdot$\, \insertshortdeck}\hfill%
  \usebeamercolor[fg]{page number in head/foot}%
  \usebeamerfont{page number in head/foot}%
  \insertframenumber\,/\,\inserttotalframenumber\kern1em\vskip2pt%
}

% ---------------------------------------------------------------------------
% Chinese text under pdflatex (book titles etc.)
\newcommand{\zh}[1]{\begin{CJK*}{UTF8}{gbsn}#1\end{CJK*}}

% ---------------------------------------------------------------------------
% Deck title block
%   \decktitle{session}{deck id}{title}{date}
\newcommand{\insertshortdeck}{}
\newcommand{\decksession}{}
\newcommand{\deckid}{}
\newcommand{\decktitletext}{}
\newcommand{\deckdate}{}
\newcommand{\decktitle}[4]{%
  \renewcommand{\decksession}{#1}%
  \renewcommand{\deckid}{#2}%
  \renewcommand{\decktitletext}{#3}%
  \renewcommand{\deckdate}{#4}%
  \renewcommand{\insertshortdeck}{Session #1 \,$\cdot$\, Deck #1.#2}%
  \title{#3}%
  \author{Zhigang Feng}%
  \date{#4}%
}
\newcommand{\makedecktitle}{%
  \begin{frame}[plain,noframenumbering]
    \begin{tikzpicture}[remember picture,overlay]
      \fill[UCRBlue] (current page.north west) rectangle ([yshift=-0.55cm]current page.north east);
      \fill[UCRGold] ([yshift=-0.55cm]current page.north west) rectangle ([yshift=-0.62cm]current page.north east);
      \node[anchor=west, text=white, font=\small\bfseries] at ([xshift=5mm,yshift=-0.275cm]current page.north west)
        {ECON 282E \,$\cdot$\, Foundations of Macroeconomics};
      \node[anchor=east, text=white, font=\small] at ([xshift=-5mm,yshift=-0.275cm]current page.north east)
        {UC Riverside \,$\cdot$\, Fall 2026};
    \end{tikzpicture}
    \vspace*{1.1cm}
    {\usebeamerfont{title}\color{black!55}\large Session \decksession\ \,$\cdot$\, Deck \decksession.\deckid\par}
    \vspace{0.25cm}
    {\usebeamerfont{title}\color{RedTitle}\LARGE\bfseries \decktitletext\par}
    \vspace{0.7cm}
    {\large Zhigang Feng\par}
    \vspace{0.15cm}
    {\color{black!65}\deckdate\par}
    \vfill
    {\footnotesize\itshape\color{black!55}Build it on paper $\;\to\;$ solve it on the machine $\;\to\;$ push it to the frontier with AI\par}
  \end{frame}}

% ---------------------------------------------------------------------------
% Section divider (plain frame, not counted)
\newcommand{\sectiondivider}[2]{%
\begin{frame}[plain,noframenumbering]
\begin{center}
\vfill
{\Large\textbf{#1}\par}
\vspace{0.35cm}
{\normalsize #2\par}
\vfill
\end{center}
\end{frame}
}

% ---------------------------------------------------------------------------
% Boxes
\newtcolorbox{conceptbox}[1][]{colback=blue!3, colframe=RedTitle!55, boxrule=0.35mm,
  arc=1mm, left=2mm, right=2mm, top=1mm, bottom=1mm, #1}
\newtcolorbox{cautionbox}[1][]{colback=red!3, colframe=red!50!black, boxrule=0.35mm,
  arc=1mm, left=2mm, right=2mm, top=1mm, bottom=1mm, #1}
\newtcolorbox{examplebox}[1][]{colback=green!3, colframe=green!45!black, boxrule=0.35mm,
  arc=1mm, left=2mm, right=2mm, top=1mm, bottom=1mm, #1}
\newtcolorbox{takeawaybox}[1][]{colback=VeryLightGray, colframe=RedTitle!55, boxrule=0.35mm,
  arc=1mm, left=2mm, right=2mm, top=1mm, bottom=1mm, #1}
\newtcolorbox{researchbox}[1][]{enhanced, colback=violet!4, colframe=violet!60!black,
  boxrule=0.35mm, arc=1mm, left=2mm, right=2mm, top=1mm, bottom=1mm,
  fonttitle=\bfseries\small, title={Research in the wild}, #1}

\newcommand{\compacttables}{\renewcommand{\arraystretch}{1.15}}
\newcommand{\src}[1]{{\tiny\color{black!55}#1}}

% ---------------------------------------------------------------------------
%% Course map: the recurring "you are here" slide for ECON 282E.
%% Loaded by course_preamble.tex. Pattern follows AI-ECON-2026 lec01_map.tex, but the map
%% shows the whole ten-session course rather than one lecture.
%%
%%   \CourseMap{s}{label}   s = 1..10 highlights session s and prints "you are here: label"
%%                          (e.g. \CourseMap{1}{1.A2}); earlier sessions get a check mark.
%%   \CourseMap{0}{}        full overview, nothing highlighted.
%%
%% Note: TikZ option lists are comma-split before TeX conditionals run, so every \ifnum
%% below sits outside [...] option lists.

\newcommand{\CourseMap}[2]{%
\begin{center}
\begin{tikzpicture}[
    sess/.style={rectangle, rounded corners=2pt, draw=black!45, minimum width=1.34cm,
        minimum height=1.12cm, text width=1.26cm, align=center, inner sep=1pt},
    stage/.style={font=\tiny\bfseries, text=black!70},
]
  \foreach \i/\d/\t/\c in {%
      1/Sep 24/Intro \\ Growth/StageTrunk,
      2/Oct 1/HA $\cdot$ OLG \\ Policy/StageBranch,
      3/Oct 8/Num.\ DP \\ Python/StageMachine,
      4/Oct 15/Numerics \\ I/StageMachine,
      5/Oct 22/Numerics \\ II/StageMachine,
      6/Oct 29/The wall \\ \textbf{Midterm}/StageWall,
      7/Nov 5/ML \\ Deep nets/StageTools,
      8/Nov 12/RL \\ HA $+$ DL/StageTools,
      9/Nov 19/Text \\ LLMs/StageData,
      10/Dec 3/Agents \\ \textbf{Projects}/StageAgents} {
    \node[sess, fill=\c] (n\i) at ({1.46*(\i-1)}, 0)
      {{\scriptsize\bfseries S\i}\\[-1pt]{\tiny\color{black!60}\d}\\[1pt]{\tiny \t}};
    \ifnum\i<#1
      \node[font=\scriptsize, text=green!45!black] at ([xshift=-2.5mm,yshift=-2.2mm]n\i.north east) {$\checkmark$};
    \fi
    \ifnum\i=#1
      \draw[RedTitle, very thick, rounded corners=2pt]
        ([xshift=-0.6mm,yshift=-0.6mm]n\i.south west) rectangle ([xshift=0.6mm,yshift=0.6mm]n\i.north east);
      % keep the label inside the picture at both ends of the row
      \ifnum\i<3
        \node[font=\scriptsize\bfseries, text=RedTitle, anchor=south west] at ([yshift=1.2mm]n\i.north west)
          {$\blacktriangledown$\,you are here: #2};
      \else\ifnum\i>8
        \node[font=\scriptsize\bfseries, text=RedTitle, anchor=south east] at ([yshift=1.2mm]n\i.north east)
          {you are here: #2\,$\blacktriangledown$};
      \else
        \node[font=\scriptsize\bfseries, text=RedTitle, anchor=south] at ([yshift=1.2mm]n\i.north)
          {you are here: #2\,$\blacktriangledown$};
      \fi\fi
    \fi
  }
  % stage band under the sessions
  \foreach \a/\b/\lab/\c in {1/1/trunk/StageTrunk, 2/2/branches/StageBranch,
      3/5/the machine $\cdot$ classical methods/StageMachine, 6/6/the wall/StageWall,
      7/8/new tools/StageTools, 9/9/new data/StageData, 10/10/colleagues/StageAgents} {
    \fill[\c] ([yshift=-1.5mm]n\a.south west) rectangle ([yshift=-4.6mm]n\b.south east);
    \path ([yshift=-1.5mm]n\a.south west) -- ([yshift=-4.6mm]n\b.south east)
      node[midway, stage] {\lab};
  }
  \node[font=\footnotesize\itshape, text=black!70] at ([yshift=-9.5mm]$(n1.south)!0.5!(n10.south)$)
    {Build it on paper $\;\to\;$ solve it on the machine $\;\to\;$ push it to the frontier with AI};
\end{tikzpicture}
\end{center}
}


\graphicspath{{./figures/}}

\lstset{language=Python, basicstyle=\ttfamily\scriptsize, frame=none,
        backgroundcolor=\color{gray!8}, aboveskip=2pt, belowskip=2pt}

\decktitle{5}{B4}{Finite Elements}{Thursday, October 22, 2026}

\begin{document}

\makedecktitle

%=============================================================================
\begin{frame}{Where We Are}
\CourseMap{5}{5.B4}
\vspace{-1mm}
\begin{center}
\small \textbf{Block B:} 5.B1 the residual $\;\cdot\;$ 5.B2 choosing a basis $\;\cdot\;$ 5.B3 Chebyshev $\;\cdot\;$ \textbf{5.B4} finite elements $\;\cdot\;$ 5.B5 determining the coefficients $\;\cdot\;$ 5.B6 the worked example.
\end{center}
\end{frame}

%=============================================================================
\begin{frame}{The Opposite Choice}
\begin{columns}[T]
\begin{column}{0.54\textwidth}
\small
A spectral basis is \textbf{global}: every $\psi_i$ is non-zero almost everywhere, so every coefficient affects the answer everywhere. That is what makes it converge geometrically on smooth functions --- and what makes a single kink contaminate the whole domain.
\medskip

A \textbf{finite element} basis is local. Chop the state space into elements, put a simple function on each, and let each basis function be non-zero only on its own small neighbourhood.
\medskip

\textbf{Then a kink is not a problem; it is a place to put an element boundary.}
\end{column}
\begin{column}{0.43\textwidth}
\begin{conceptbox}[title=\textbf{Where FEM comes from}]\small
Engineering, for PDEs on awkward domains --- irregular boundaries, corners, complex geometry. The mesh can be shaped to fit the problem.
\end{conceptbox}
\medskip
\begin{examplebox}\footnotesize
In macroeconomics the ``irregular boundary'' is usually a constraint: $a'\ge-\underline{a}$, or $i_t\ge0$. Put a node exactly there.
\end{examplebox}
\end{column}
\end{columns}
\vspace{1mm}
\src{Quant\_Macro Lec.~9, ``Finite Element Method (FEM)''.}
\end{frame}

%=============================================================================
\begin{frame}{The Recipe}
\begin{columns}[T]
\begin{column}{0.56\textwidth}
\small
\begin{enumerate}\small\setlength{\itemsep}{1.2mm}
  \item \textbf{Partition} the state space into elements --- intervals in one dimension, triangles or rectangles in more.
  \item \textbf{Choose a basis on each element}, typically piecewise polynomial, and typically \textbf{$1$ at one node and $0$ at all the others}. Summed, they form a \textbf{partition of unity}.
  \item \textbf{Substitute} the approximant into the equilibrium conditions and integrate over each element, often by parts.
  \item \textbf{Assemble} the element equations into one global algebraic system.
  \item \textbf{Solve} for the nodal coefficients, and \textbf{paste} the pieces together so the result is continuous.
\end{enumerate}
\end{column}
\begin{column}{0.41\textwidth}
\begin{takeawaybox}\footnotesize
Steps 1--2 are the basis choice of 5.B2; steps 3--4 are the projection of 5.B5. \textbf{Nothing in the framework changed} --- only $\{\psi_i\}$.
\end{takeawaybox}
\medskip
\begin{examplebox}\footnotesize
``Paste it together to ensure continuity'' is doing real work: the elements are solved jointly, not one at a time.
\end{examplebox}
\end{column}
\end{columns}
\vspace{1mm}
\src{Quant\_Macro Lec.~9, ``Finite Element Method (FEM)'' and ``Structure''.}
\end{frame}

%=============================================================================
\begin{frame}{The Simplest Basis: Tent Functions}
\begin{columns}[T]
\begin{column}{0.5\textwidth}
\small
Elements are small, so a linear basis is usually enough:
\[
g_{i}(x)=\begin{cases}
\dfrac{x-x_{i-1}}{x_{i}-x_{i-1}} & x\in[x_{i-1},x_{i}]\\[2mm]
\dfrac{x_{i+1}-x}{x_{i+1}-x_{i}} & x\in[x_{i},x_{i+1}]\\[1mm]
0 & \text{elsewhere}
\end{cases}
\]
Each $g_i$ is $1$ at its own node, $0$ at every other node, and non-zero only on $[x_{i-1},x_{i+1}]$.
\medskip

Then $d^{n}(x)=\sum_i\theta_i g_i(x)$ and $\theta_i$ is simply \textbf{the value of the policy at node $i$} --- the coefficients are interpretable in a way Chebyshev coefficients never are.
\end{column}
\begin{column}{0.47\textwidth}
\begin{center}
\begin{tikzpicture}
\begin{axis}[
    width=6.9cm, height=4.3cm,
    xmin=0, xmax=1, ymin=-0.08, ymax=1.25,
    xlabel={\scriptsize $x$},
    tick label style={font=\scriptsize}, label style={font=\scriptsize},
    scaled ticks=false, /pgf/number format/fixed,
    /pgf/number format/precision=2,
    axis x line*=bottom, axis y line*=left,
    legend style={font=\tiny, draw=none, fill=none,
                  at={(0.5,-0.30)}, anchor=north, legend columns=2}]
  \addplot[very thick, black!40] coordinates {
      (0.0000,1.0000) (0.0050,1.0000) (0.0100,1.0000) (0.0150,1.0000) (0.0200,1.0000) 
      (0.0250,1.0000) (0.0300,1.0000) (0.0350,1.0000) (0.0400,1.0000) (0.0450,1.0000) 
      (0.0500,1.0000) (0.0550,1.0000) (0.0600,1.0000) (0.0650,1.0000) (0.0700,1.0000) 
      (0.0750,1.0000) (0.0800,1.0000) (0.0850,1.0000) (0.0900,1.0000) (0.0950,1.0000) 
      (0.1000,1.0000) (0.1050,1.0000) (0.1100,1.0000) (0.1150,1.0000) (0.1200,1.0000) 
      (0.1250,1.0000) (0.1300,1.0000) (0.1350,1.0000) (0.1400,1.0000) (0.1450,1.0000) 
      (0.1500,1.0000) (0.1550,1.0000) (0.1600,1.0000) (0.1650,1.0000) (0.1700,1.0000) 
      (0.1750,1.0000) (0.1800,1.0000) (0.1850,1.0000) (0.1900,1.0000) (0.1950,1.0000) 
      (0.2000,1.0000) (0.2050,1.0000) (0.2100,1.0000) (0.2150,1.0000) (0.2200,1.0000) 
      (0.2250,1.0000) (0.2300,1.0000) (0.2350,1.0000) (0.2400,1.0000) (0.2450,1.0000) 
      (0.2500,1.0000) (0.2550,1.0000) (0.2600,1.0000) (0.2650,1.0000) (0.2700,1.0000) 
      (0.2750,1.0000) (0.2800,1.0000) (0.2850,1.0000) (0.2900,1.0000) (0.2950,1.0000) 
      (0.3000,1.0000) (0.3050,1.0000) (0.3100,1.0000) (0.3150,1.0000) (0.3200,1.0000) 
      (0.3250,1.0000) (0.3300,1.0000) (0.3350,1.0000) (0.3400,1.0000) (0.3450,1.0000) 
      (0.3500,1.0000) (0.3550,1.0000) (0.3600,1.0000) (0.3650,1.0000) (0.3700,1.0000) 
      (0.3750,1.0000) (0.3800,1.0000) (0.3850,1.0000) (0.3900,1.0000) (0.3950,1.0000) 
      (0.4000,1.0000) (0.4050,1.0000) (0.4100,1.0000) (0.4150,1.0000) (0.4200,1.0000) 
      (0.4250,1.0000) (0.4300,1.0000) (0.4350,1.0000) (0.4400,1.0000) (0.4450,1.0000) 
      (0.4500,1.0000) (0.4550,1.0000) (0.4600,1.0000) (0.4650,1.0000) (0.4700,1.0000) 
      (0.4750,1.0000) (0.4800,1.0000) (0.4850,1.0000) (0.4900,1.0000) (0.4950,1.0000) 
      (0.5000,1.0000) (0.5050,1.0000) (0.5100,1.0000) (0.5150,1.0000) (0.5200,1.0000) 
      (0.5250,1.0000) (0.5300,1.0000) (0.5350,1.0000) (0.5400,1.0000) (0.5450,1.0000) 
      (0.5500,1.0000) (0.5550,1.0000) (0.5600,1.0000) (0.5650,1.0000) (0.5700,1.0000) 
      (0.5750,1.0000) (0.5800,1.0000) (0.5850,1.0000) (0.5900,1.0000) (0.5950,1.0000) 
      (0.6000,1.0000) (0.6050,1.0000) (0.6100,1.0000) (0.6150,1.0000) (0.6200,1.0000) 
      (0.6250,1.0000) (0.6300,1.0000) (0.6350,1.0000) (0.6400,1.0000) (0.6450,1.0000) 
      (0.6500,1.0000) (0.6550,1.0000) (0.6600,1.0000) (0.6650,1.0000) (0.6700,1.0000) 
      (0.6750,1.0000) (0.6800,1.0000) (0.6850,1.0000) (0.6900,1.0000) (0.6950,1.0000) 
      (0.7000,1.0000) (0.7050,1.0000) (0.7100,1.0000) (0.7150,1.0000) (0.7200,1.0000) 
      (0.7250,1.0000) (0.7300,1.0000) (0.7350,1.0000) (0.7400,1.0000) (0.7450,1.0000) 
      (0.7500,1.0000) (0.7550,1.0000) (0.7600,1.0000) (0.7650,1.0000) (0.7700,1.0000) 
      (0.7750,1.0000) (0.7800,1.0000) (0.7850,1.0000) (0.7900,1.0000) (0.7950,1.0000) 
      (0.8000,1.0000) (0.8050,1.0000) (0.8100,1.0000) (0.8150,1.0000) (0.8200,1.0000) 
      (0.8250,1.0000) (0.8300,1.0000) (0.8350,1.0000) (0.8400,1.0000) (0.8450,1.0000) 
      (0.8500,1.0000) (0.8550,1.0000) (0.8600,1.0000) (0.8650,1.0000) (0.8700,1.0000) 
      (0.8750,1.0000) (0.8800,1.0000) (0.8850,1.0000) (0.8900,1.0000) (0.8950,1.0000) 
      (0.9000,1.0000) (0.9050,1.0000) (0.9100,1.0000) (0.9150,1.0000) (0.9200,1.0000) 
      (0.9250,1.0000) (0.9300,1.0000) (0.9350,1.0000) (0.9400,1.0000) (0.9450,1.0000) 
      (0.9500,1.0000) (0.9550,1.0000) (0.9600,1.0000) (0.9650,1.0000) (0.9700,1.0000) 
      (0.9750,1.0000) (0.9800,1.0000) (0.9850,1.0000) (0.9900,1.0000) (0.9950,1.0000) 
      (1.0000,1.0000) };
  \addlegendentry{$\sum_i g_i \equiv 1$}
  \addplot[very thick, RedTitle] coordinates {
      (0.0000,0.0000) (0.0050,0.0000) (0.0100,0.0000) (0.0150,0.0000) (0.0200,0.0000) 
      (0.0250,0.0000) (0.0300,0.0000) (0.0350,0.0000) (0.0400,0.0000) (0.0450,0.0000) 
      (0.0500,0.0000) (0.0550,0.0000) (0.0600,0.0000) (0.0650,0.0000) (0.0700,0.0000) 
      (0.0750,0.0000) (0.0800,0.0000) (0.0850,0.0000) (0.0900,0.0000) (0.0950,0.0000) 
      (0.1000,0.0000) (0.1050,0.0000) (0.1100,0.0000) (0.1150,0.0000) (0.1200,0.0000) 
      (0.1250,0.0000) (0.1300,0.0000) (0.1350,0.0000) (0.1400,0.0000) (0.1450,0.0000) 
      (0.1500,0.0000) (0.1550,0.0000) (0.1600,0.0000) (0.1650,0.0000) (0.1700,0.0000) 
      (0.1750,0.0000) (0.1800,0.0000) (0.1850,0.0000) (0.1900,0.0000) (0.1950,0.0000) 
      (0.2000,0.0000) (0.2050,0.0000) (0.2100,0.0000) (0.2150,0.0000) (0.2200,0.0000) 
      (0.2250,0.0000) (0.2300,0.0000) (0.2350,0.0000) (0.2400,0.0000) (0.2450,0.0000) 
      (0.2500,0.0000) (0.2550,0.0000) (0.2600,0.0000) (0.2650,0.0000) (0.2700,0.0000) 
      (0.2750,0.0000) (0.2800,0.0000) (0.2850,0.0000) (0.2900,0.0000) (0.2950,0.0000) 
      (0.3000,0.0000) (0.3050,0.0000) (0.3100,0.0000) (0.3150,0.0000) (0.3200,0.0000) 
      (0.3250,0.0000) (0.3300,0.0000) (0.3350,0.0000) (0.3400,0.0000) (0.3450,0.0000) 
      (0.3500,0.0000) (0.3550,0.0000) (0.3600,0.0000) (0.3650,0.0000) (0.3700,0.0000) 
      (0.3750,0.0000) (0.3800,0.0000) (0.3850,0.0000) (0.3900,0.0000) (0.3950,0.0000) 
      (0.4000,0.0000) (0.4050,0.0000) (0.4100,0.0000) (0.4150,0.0000) (0.4200,0.0000) 
      (0.4250,0.0000) (0.4300,0.0000) (0.4350,0.0000) (0.4400,0.0000) (0.4450,0.0000) 
      (0.4500,0.0000) (0.4550,0.0000) (0.4600,0.0000) (0.4650,0.0000) (0.4700,0.0000) 
      (0.4750,0.0000) (0.4800,0.0000) (0.4850,0.0000) (0.4900,0.0000) (0.4950,0.0000) 
      (0.5000,0.0000) (0.5050,0.0000) (0.5100,0.0000) (0.5150,0.0000) (0.5200,0.0000) 
      (0.5250,0.0000) (0.5300,0.0000) (0.5350,0.0000) (0.5400,0.0000) (0.5450,0.0000) 
      (0.5500,0.0000) (0.5550,0.0000) (0.5600,0.0000) (0.5650,0.0000) (0.5700,0.0000) 
      (0.5750,0.0000) (0.5800,0.0000) (0.5850,0.0000) (0.5900,0.0000) (0.5950,0.0000) 
      (0.6000,0.0000) (0.6050,0.0000) (0.6100,0.0000) (0.6150,0.0000) (0.6200,0.0000) 
      (0.6250,0.0000) (0.6300,0.0000) (0.6350,0.0000) (0.6400,0.0000) (0.6450,0.0000) 
      (0.6500,0.0000) (0.6550,0.0000) (0.6600,0.0000) (0.6650,0.0000) (0.6700,0.0000) 
      (0.6750,0.0000) (0.6800,0.0000) (0.6850,0.0000) (0.6900,0.0000) (0.6950,0.0000) 
      (0.7000,0.0000) (0.7050,0.0000) (0.7100,0.0000) (0.7150,0.0000) (0.7200,0.0000) 
      (0.7250,0.0000) (0.7300,0.0000) (0.7350,0.0000) (0.7400,0.0000) (0.7450,0.0000) 
      (0.7500,0.0000) (0.7550,0.0000) (0.7600,0.0000) (0.7650,0.0000) (0.7700,0.0000) 
      (0.7750,0.0000) (0.7800,0.0000) (0.7850,0.0000) (0.7900,0.0000) (0.7950,0.0000) 
      (0.8000,0.0000) (0.8050,0.0000) (0.8100,0.0000) (0.8150,0.0000) (0.8200,0.0000) 
      (0.8250,0.0000) (0.8300,0.0000) (0.8350,0.0000) (0.8400,0.0000) (0.8450,0.0000) 
      (0.8500,0.0000) (0.8550,0.0000) (0.8600,0.0000) (0.8650,0.0000) (0.8700,0.0000) 
      (0.8750,0.0000) (0.8800,0.0000) (0.8850,0.0000) (0.8900,0.0000) (0.8950,0.0000) 
      (0.9000,0.0000) (0.9050,0.0000) (0.9100,0.0000) (0.9150,0.0000) (0.9200,0.0000) 
      (0.9250,0.0000) (0.9300,0.0000) (0.9350,0.0000) (0.9400,0.0000) (0.9450,0.0000) 
      (0.9500,0.0000) (0.9550,0.0000) (0.9600,0.0000) (0.9650,0.0000) (0.9700,0.0000) 
      (0.9750,0.0000) (0.9800,0.0000) (0.9850,0.0000) (0.9900,0.0000) (0.9950,0.0000) 
      (1.0000,0.0000) };
  \addlegendentry{$g_3$}
  \addplot[very thick, orange!85!black] coordinates {
      (0.0000,0.0000) (0.0050,0.0000) (0.0100,0.0000) (0.0150,0.0000) (0.0200,0.0000) 
      (0.0250,0.0000) (0.0300,0.0000) (0.0350,0.0000) (0.0400,0.0000) (0.0450,0.0000) 
      (0.0500,0.0000) (0.0550,0.0000) (0.0600,0.0000) (0.0650,0.0000) (0.0700,0.0000) 
      (0.0750,0.0000) (0.0800,0.0000) (0.0850,0.0000) (0.0900,0.0000) (0.0950,0.0000) 
      (0.1000,0.0000) (0.1050,0.0000) (0.1100,0.0000) (0.1150,0.0000) (0.1200,0.0000) 
      (0.1250,0.0000) (0.1300,0.0000) (0.1350,0.0000) (0.1400,0.0000) (0.1450,0.0000) 
      (0.1500,0.0000) (0.1550,0.0000) (0.1600,0.0000) (0.1650,0.0000) (0.1700,0.0000) 
      (0.1750,0.0000) (0.1800,0.0000) (0.1850,0.0000) (0.1900,0.0000) (0.1950,0.0000) 
      (0.2000,0.0000) (0.2050,0.0000) (0.2100,0.0000) (0.2150,0.0000) (0.2200,0.0000) 
      (0.2250,0.0000) (0.2300,0.0000) (0.2350,0.0000) (0.2400,0.0000) (0.2450,0.0000) 
      (0.2500,0.0000) (0.2550,0.0000) (0.2600,0.0000) (0.2650,0.0000) (0.2700,0.0000) 
      (0.2750,0.0000) (0.2800,0.0000) (0.2850,0.0000) (0.2900,0.0000) (0.2950,0.0000) 
      (0.3000,0.0000) (0.3050,0.0000) (0.3100,0.0000) (0.3150,0.0000) (0.3200,0.0000) 
      (0.3250,0.0000) (0.3300,0.0000) (0.3350,0.0000) (0.3400,0.0000) (0.3450,0.0000) 
      (0.3500,0.0000) (0.3550,0.0000) (0.3600,0.0000) (0.3650,0.0000) (0.3700,0.0000) 
      (0.3750,0.0000) (0.3800,0.0000) (0.3850,0.0000) (0.3900,0.0000) (0.3950,0.0000) 
      (0.4000,0.0000) (0.4050,0.0000) (0.4100,0.0000) (0.4150,0.0000) (0.4200,0.0000) 
      (0.4250,0.0000) (0.4300,0.0000) (0.4350,0.0000) (0.4400,0.0000) (0.4450,0.0000) 
      (0.4500,0.0000) (0.4550,0.0000) (0.4600,0.0000) (0.4650,0.0000) (0.4700,0.0000) 
      (0.4750,0.0000) (0.4800,0.0000) (0.4850,0.0000) (0.4900,0.0000) (0.4950,0.0000) 
      (0.5000,0.0000) (0.5050,0.0000) (0.5100,0.0000) (0.5150,0.0000) (0.5200,0.0000) 
      (0.5250,0.0000) (0.5300,0.0000) (0.5350,0.0000) (0.5400,0.0000) (0.5450,0.0000) 
      (0.5500,0.0000) (0.5550,0.0000) (0.5600,0.0000) (0.5650,0.0000) (0.5700,0.0000) 
      (0.5750,0.0000) (0.5800,0.0000) (0.5850,0.0000) (0.5900,0.0000) (0.5950,0.0000) 
      (0.6000,0.0000) (0.6050,0.0000) (0.6100,0.0000) (0.6150,0.0000) (0.6200,0.0000) 
      (0.6250,0.0000) (0.6300,0.0000) (0.6350,0.0000) (0.6400,0.0000) (0.6450,0.0000) 
      (0.6500,0.0000) (0.6550,0.0000) (0.6600,0.0000) (0.6650,0.0000) (0.6700,0.0000) 
      (0.6750,0.0000) (0.6800,0.0000) (0.6850,0.0000) (0.6900,0.0000) (0.6950,0.0000) 
      (0.7000,0.0000) (0.7050,0.0000) (0.7100,0.0000) (0.7150,0.0000) (0.7200,0.0000) 
      (0.7250,0.0000) (0.7300,0.0000) (0.7350,0.0000) (0.7400,0.0000) (0.7450,0.0000) 
      (0.7500,0.0000) (0.7550,0.0000) (0.7600,0.0000) (0.7650,0.0000) (0.7700,0.0000) 
      (0.7750,0.0000) (0.7800,0.0000) (0.7850,0.0000) (0.7900,0.0000) (0.7950,0.0000) 
      (0.8000,0.0000) (0.8050,0.0000) (0.8100,0.0000) (0.8150,0.0000) (0.8200,0.0000) 
      (0.8250,0.0000) (0.8300,0.0000) (0.8350,0.0000) (0.8400,0.0000) (0.8450,0.0000) 
      (0.8500,0.0000) (0.8550,0.0000) (0.8600,0.0000) (0.8650,0.0000) (0.8700,0.0000) 
      (0.8750,0.0000) (0.8800,0.0000) (0.8850,0.0000) (0.8900,0.0000) (0.8950,0.0000) 
      (0.9000,0.0000) (0.9050,0.0000) (0.9100,0.0000) (0.9150,0.0000) (0.9200,0.0000) 
      (0.9250,0.0000) (0.9300,0.0000) (0.9350,0.0000) (0.9400,0.0000) (0.9450,0.0000) 
      (0.9500,0.0000) (0.9550,0.0000) (0.9600,0.0000) (0.9650,0.0000) (0.9700,0.0000) 
      (0.9750,0.0000) (0.9800,0.0000) (0.9850,0.0000) (0.9900,0.0000) (0.9950,0.0000) 
      (1.0000,0.0000) };
  \addlegendentry{$g_4$}
  \addplot[very thick, black!65] coordinates {
      (0.0000,0.0000) (0.0050,0.0000) (0.0100,0.0000) (0.0150,0.0000) (0.0200,0.0000) 
      (0.0250,0.0000) (0.0300,0.0000) (0.0350,0.0000) (0.0400,0.0000) (0.0450,0.0000) 
      (0.0500,0.0000) (0.0550,0.0000) (0.0600,0.0000) (0.0650,0.0000) (0.0700,0.0000) 
      (0.0750,0.0000) (0.0800,0.0000) (0.0850,0.0000) (0.0900,0.0000) (0.0950,0.0000) 
      (0.1000,0.0000) (0.1050,0.0000) (0.1100,0.0000) (0.1150,0.0000) (0.1200,0.0000) 
      (0.1250,0.0000) (0.1300,0.0000) (0.1350,0.0000) (0.1400,0.0000) (0.1450,0.0000) 
      (0.1500,0.0000) (0.1550,0.0000) (0.1600,0.0000) (0.1650,0.0000) (0.1700,0.0000) 
      (0.1750,0.0000) (0.1800,0.0000) (0.1850,0.0000) (0.1900,0.0000) (0.1950,0.0000) 
      (0.2000,0.0000) (0.2050,0.0000) (0.2100,0.0000) (0.2150,0.0000) (0.2200,0.0000) 
      (0.2250,0.0000) (0.2300,0.0000) (0.2350,0.0000) (0.2400,0.0000) (0.2450,0.0000) 
      (0.2500,0.0000) (0.2550,0.0000) (0.2600,0.0000) (0.2650,0.0000) (0.2700,0.0000) 
      (0.2750,0.0000) (0.2800,0.0000) (0.2850,0.0000) (0.2900,0.0000) (0.2950,0.0000) 
      (0.3000,0.0000) (0.3050,0.0000) (0.3100,0.0000) (0.3150,0.0000) (0.3200,0.0000) 
      (0.3250,0.0000) (0.3300,0.0000) (0.3350,0.0000) (0.3400,0.0000) (0.3450,0.0000) 
      (0.3500,0.0000) (0.3550,0.0000) (0.3600,0.0000) (0.3650,0.0000) (0.3700,0.0000) 
      (0.3750,0.0000) (0.3800,0.0000) (0.3850,0.0000) (0.3900,0.0000) (0.3950,0.0000) 
      (0.4000,0.0000) (0.4050,0.0000) (0.4100,0.0000) (0.4150,0.0000) (0.4200,0.0000) 
      (0.4250,0.0000) (0.4300,0.0000) (0.4350,0.0000) (0.4400,0.0000) (0.4450,0.0000) 
      (0.4500,0.0000) (0.4550,0.0000) (0.4600,0.0000) (0.4650,0.0000) (0.4700,0.0000) 
      (0.4750,0.0000) (0.4800,0.0000) (0.4850,0.0000) (0.4900,0.0000) (0.4950,0.0000) 
      (0.5000,0.0000) (0.5050,0.0000) (0.5100,0.0000) (0.5150,0.0000) (0.5200,0.0000) 
      (0.5250,0.0000) (0.5300,0.0000) (0.5350,0.0000) (0.5400,0.0000) (0.5450,0.0000) 
      (0.5500,0.0000) (0.5550,0.0000) (0.5600,0.0000) (0.5650,0.0000) (0.5700,0.0000) 
      (0.5750,0.0000) (0.5800,0.0000) (0.5850,0.0000) (0.5900,0.0000) (0.5950,0.0000) 
      (0.6000,0.0000) (0.6050,0.0000) (0.6100,0.0000) (0.6150,0.0000) (0.6200,0.0000) 
      (0.6250,0.0000) (0.6300,0.0000) (0.6350,0.0000) (0.6400,0.0000) (0.6450,0.0000) 
      (0.6500,0.0000) (0.6550,0.0000) (0.6600,0.0000) (0.6650,0.0000) (0.6700,0.0000) 
      (0.6750,0.0000) (0.6800,0.0000) (0.6850,0.0000) (0.6900,0.0000) (0.6950,0.0000) 
      (0.7000,0.0000) (0.7050,0.0000) (0.7100,0.0000) (0.7150,0.0000) (0.7200,0.0000) 
      (0.7250,0.0000) (0.7300,0.0000) (0.7350,0.0000) (0.7400,0.0000) (0.7450,0.0000) 
      (0.7500,0.0000) (0.7550,0.0000) (0.7600,0.0000) (0.7650,0.0000) (0.7700,0.0000) 
      (0.7750,0.0000) (0.7800,0.0000) (0.7850,0.0000) (0.7900,0.0000) (0.7950,0.0000) 
      (0.8000,0.0000) (0.8050,0.0000) (0.8100,0.0000) (0.8150,0.0000) (0.8200,0.0000) 
      (0.8250,0.0000) (0.8300,0.0000) (0.8350,0.0000) (0.8400,0.0000) (0.8450,0.0000) 
      (0.8500,0.0000) (0.8550,0.0000) (0.8600,0.0000) (0.8650,0.0000) (0.8700,0.0000) 
      (0.8750,0.0000) (0.8800,0.0000) (0.8850,0.0000) (0.8900,0.0000) (0.8950,0.0000) 
      (0.9000,0.0000) (0.9050,0.0000) (0.9100,0.0000) (0.9150,0.0000) (0.9200,0.0000) 
      (0.9250,0.0000) (0.9300,0.0000) (0.9350,0.0000) (0.9400,0.0000) (0.9450,0.0000) 
      (0.9500,0.0000) (0.9550,0.0000) (0.9600,0.0000) (0.9650,0.0000) (0.9700,0.0000) 
      (0.9750,0.0000) (0.9800,0.0000) (0.9850,0.0000) (0.9900,0.0000) (0.9950,0.0000) 
      (1.0000,0.0000) };
  \addlegendentry{$g_5$}
\end{axis}
\end{tikzpicture}
\end{center}
\end{column}
\end{columns}
\vspace{1mm}
\src{Quant\_Macro Lec.~9, ``Structure''. Nine nodes on $[0,1]$; measured $\max_x|\sum_i g_i(x)-1| = 0$.}
\end{frame}

%=============================================================================
\begin{frame}{Check the Partition of Unity Before You Solve Anything}
\begin{columns}[T]
\begin{column}{0.53\textwidth}
\small
The tent function has two branches and it is easy to get the second one's sign wrong. Write
\[
  \frac{x_{i-1}-x}{x_{i}-x_{i-1}}\qquad\text{for }x\in[x_i,x_{i+1}]
\]
instead of $(x_{i+1}-x)/(x_{i+1}-x_i)$, and evaluate just to the right of the node:
\[
  g_i(x_i^{+})=\frac{x_{i-1}-x_i}{x_i-x_{i-1}}=\mathbf{-1}.
\]
A basis function that is $+1$ arriving at its node from the left and $-1$ leaving it to the right is
discontinuous and negative. Summed over nine nodes it is off by as much as $\mathbf{2}$; the correct
basis is off by $\mathbf{0}$.
\end{column}
\begin{column}{0.44\textwidth}
\begin{cautionbox}[title=\textbf{Why it matters}]\small
The partition-of-unity property is what makes $\theta_i$ the policy value at node $i$, and what makes
the assembled approximation continuous. Lose it and every interpretation on the previous slide fails
--- silently, because the solver still returns coefficients.
\end{cautionbox}
\medskip
\begin{takeawaybox}\footnotesize
\textbf{One line of code separates the two}: sum the basis over a fine grid and compare to $1$. Do it
before solving anything, on every basis you write yourself.
\end{takeawaybox}
\end{column}
\end{columns}
\vspace{1mm}
\src{Ledger \texttt{fem}: $\max_x|\sum_i g_i(x)-1| = 0$ for the correct basis and $2.0$ with the sign flipped; value just past the node $=-1$.}
\end{frame}

%=============================================================================
\begin{frame}{Why the System Is Sparse}
\begin{columns}[T]
\begin{column}{0.52\textwidth}
\small
Basis function $g_i$ is non-zero only on $[x_{i-1},x_{i+1}]$. So in any integral of the form $\int g_i\,g_j$ or $\int g_i R$, the integrand is zero unless $g_i$ and $g_j$ \textbf{share an element}.
\medskip

The assembled system $A\theta=b$ therefore has non-zeros only on the diagonal and its immediate neighbours: it is \textbf{banded}, with bandwidth $1$.
\medskip

Measured, on nine nodes: \textbf{25 non-zero entries out of 81}, a density of $30.9\%$ --- and the density falls like $3/n$ as the mesh is refined. At $500$ nodes it is under $1\%$.
\end{column}
\begin{column}{0.45\textwidth}
\begin{conceptbox}[title=\textbf{What sparsity buys}]\small
A dense solve is $O(n^{3})$; a banded one is $O(n)$. That is the trade for FEM's slow convergence: many more coefficients, each one far cheaper.
\end{conceptbox}
\medskip
\begin{examplebox}\footnotesize
And you never actually invert it --- you factor it, which for a banded matrix is linear in the number of nodes.
\end{examplebox}
\end{column}
\end{columns}
\vspace{1mm}
\src{Quant\_Macro Lec.~9, ``Why Does This Lead to a Sparse Matrix?''. Ledger \texttt{fem}: $25/81$ non-zeros, bandwidth $1$.}
\end{frame}

%=============================================================================
\begin{frame}{Three Ways to Refine}
\small
\begin{center}
\begin{tabularx}{0.97\linewidth}{@{}>{\raggedright\arraybackslash}p{2.3cm} >{\raggedright\arraybackslash}X@{}}
\toprule
\textbf{h-refinement} & Subdivide every element, improving resolution \emph{uniformly} over the domain. The obvious move, and usually the wasteful one. \\
\textbf{r-refinement} & Subdivide only where the solution is strongly nonlinear. The mesh follows the economics --- more elements near a binding constraint, few where the policy is nearly linear. \\
\textbf{p-refinement} & Raise the polynomial order \emph{inside} each element. Push it far enough and you have a hybrid of finite and spectral methods: \textbf{spectral elements}. \\
\bottomrule
\end{tabularx}
\end{center}
\medskip
\begin{columns}[T]
\begin{column}{0.49\textwidth}
\begin{conceptbox}[title=\textbf{The design space is continuous}]\footnotesize
One element with a degree-$n$ basis is a spectral method. $n$ elements with a degree-$1$ basis is FEM. Everything between is available.
\end{conceptbox}
\end{column}
\begin{column}{0.47\textwidth}
\begin{examplebox}\footnotesize
r-refinement is the one that pays in macro, and it is the same instinct as the log-spaced asset grid of 3.A2: put the points where the curvature is.
\end{examplebox}
\end{column}
\end{columns}
\vspace{1mm}
\src{Quant\_Macro Lec.~9, ``Three different refinements''.}
\end{frame}

%=============================================================================
\begin{frame}{Finite Elements Against Smolyak}
\small
\begin{center}
\begin{tabularx}{0.97\linewidth}{@{}>{\raggedright\arraybackslash}p{2.0cm} >{\raggedright\arraybackslash}X >{\raggedright\arraybackslash}X@{}}
\toprule
 & \textbf{Finite elements} & \textbf{Smolyak} \\
\midrule
Basis & piecewise polynomial, \textbf{local}; the solution may differ element by element & \textbf{global} Chebyshev or splines on a sparse index set \\
Grid & a \textbf{mesh}, which can follow irregular boundaries and be adapted locally & tensor product of 1D grids, thinned by the sparse rule \\
Strength & \textbf{low to moderate dimension} with awkward domains or kinks & \textbf{moderate to high dimension} with smooth solutions \\
Weakness & heavy in high dimension --- the mesh itself is $O(n^{d})$ & less natural for irregular domains and non-smooth solutions \\
\bottomrule
\end{tabularx}
\end{center}
\medskip
\begin{center}
\begin{minipage}{0.88\linewidth}
\begin{takeawaybox}\footnotesize
They fail in \textbf{opposite} directions, which is the useful thing about the comparison: dimension kills FEM, and non-smoothness kills Smolyak. Knowing which one you face is the choice.
\end{takeawaybox}
\end{minipage}
\end{center}
\vspace{1mm}
\src{Quant\_Macro Lec.~9, ``Finite Elements vs.\ Smolyak (I)'' and ``(II)''.}
\end{frame}

%=============================================================================
\begin{frame}{Takeaways}
\begin{takeawaybox}
\begin{itemize}\small
  \item Finite elements are the \textbf{local} basis: chop the domain, put a simple function on each piece. A kink stops being a problem and becomes a place to put an element boundary.
  \item The tent basis is $1$ at its own node and $0$ at every other, so the coefficients \textbf{are} the policy values at the nodes --- directly interpretable, unlike Chebyshev coefficients.
  \item \textbf{Check the partition of unity before solving anything.} One sign slip in the tent formula gives a basis worth $\mathbf{-1}$ just past its own node, off by $2$ when summed; the correct one is exact to $0$, and one line of code tells them apart.
  \item Compact support makes the assembled system \textbf{banded}: measured $25$ non-zeros out of $81$ on nine nodes, and the density falls like $3/n$. A banded solve is $O(n)$, not $O(n^{3})$.
  \item \textbf{h-, r-, p-refinement}: more elements, better-placed elements, or higher order within elements. Push $p$ far enough and you have spectral elements --- the two families are ends of one design space.
  \item FEM and Smolyak fail in \textbf{opposite} directions: dimension defeats FEM, non-smoothness defeats Smolyak.
\end{itemize}
\end{takeawaybox}
\end{frame}

%=============================================================================
\begin{frame}{Bridge: What Does ``Small'' Mean?}
\begin{columns}[T]
\begin{column}{0.53\textwidth}
\small
Three decks have now been spent on the first of 5.B1's two questions --- which basis. Spectral or local, Chebyshev or tents, dense or sparse.
\medskip

The second question has been deferred every time it came up: having written $R(x;\theta)$, what exactly do we do to it? ``Make it small'' is not an algorithm, and the choices are not equivalent.
\end{column}
\begin{column}{0.44\textwidth}
\begin{conceptbox}[title=\textbf{Next (5.B5)}]\small
Weighted residuals, and the six answers: least squares, subdomain, moments, collocation, orthogonal collocation, and Galerkin. All six run on the same model, with the same basis --- and the results compared.
\end{conceptbox}
\end{column}
\end{columns}
\end{frame}

\end{document}
